% GENERATED FILE. Edit canonical lessons, then run npm run generate:books.
% canonical-source-hash: fnv1a64:e0b409880fbf7448
% canonical-lessons: ML-C501-mrgasala, ML-C501-udyanam, ML-C501-timimgalam, ML-C501-pravesanaphisu, ML-C501-olukiyettuka

\chapter{Zoo and Shore}
\label{ch:ml-zoo-and-shore}
\hlchaptermodality{501}

\begin{chapteropening}
\textbf{By the end of this chapter:} \emph{I can say a zoo, a public garden, a whale, an entry fee and to drift in.}

Complete the last lesson of chapter 501: I can say a zoo, a public garden, a whale, an entry fee and to drift in.
\end{chapteropening}

\section[\texorpdfstring{\ml{മൃഗശാല} (mṛgaśāla)}{മൃഗശാല (mṛgaśāla)}]{\ml{മൃഗശാല} (mṛgaśāla) — a zoo}

\label{lesson:ML-C501-mrgasala}



\begin{quote}
Before the new one: say the Malayalam for to roost, then the Malayalam for to build a nest.
\end{quote}

\subsection*{You'll want to know: \ml{മൃഗശാല}}
\textbf{\ml{മൃഗശാല}} — \emph{mṛgaśāla} — \textquotedblleft{}a zoo\textquotedblright{}.

\begin{quote}
\ml{കുട്ടികൾ} \ml{മൃഗശാലയിൽ} \ml{സിംഹത്തെ} \ml{കണ്ടു}.
\end{quote}

\begin{quote}
\emph{kuṭṭikaḷ mṛgaśālayil siṁhatte kaṇṭu.}
\end{quote}

\begin{quote}
The children saw a lion at the zoo.
\end{quote}

\ml{മൃഗം}, an animal, plus \ml{ശാല}, a hall or house.

\begin{cousinweb}
\textbf{\ml{മൃഗശാല}} (\emph{mṛgaśāla}) is a Malayalam compound built from two Sanskrit words, \emph{mṛga}, \textquotedblleft{}a wild animal\textquotedblright{}, and \emph{śālā}, \textquotedblleft{}a house, a hall\textquotedblright{} (Monier-Williams).
\end{cousinweb}

\begin{grammarlens}[title={What you've built}]
One more word for getting around.
\end{grammarlens}

\subsection*{Your turn}
\begin{itemize}
\raggedright
  \item \emph{Say it:} \emph{mṛgaśāla}
  \item \emph{Say it:} \emph{mṛgaśāla}, once more
  \item \emph{Say it:} \emph{kuṭṭikaḷ mṛgaśālayil siṁhatte kaṇṭu.} — The children saw a lion at the zoo.
  \item \emph{Recall:} say the Malayalam for to roost, then the Malayalam for to build a nest, then say \emph{mṛgaśāla} again
\end{itemize}

\begin{tcolorbox}[breakable,colback=teal!4,colframe=teal!35!black,title={Before you move on}]
What does \ml{മൃഗശാല} mean? (\textquotedblleft{}A zoo\textquotedblright{}.) Say it once more.
\end{tcolorbox}
\section[\texorpdfstring{\ml{ഉദ്യാനം} (udyānaṁ)}{ഉദ്യാനം (udyānaṁ)}]{\ml{ഉദ്യാനം} (udyānaṁ) — a public garden; a park}

\label{lesson:ML-C501-udyanam}



\begin{quote}
Before the new one: say the Malayalam for to build a nest, then the Malayalam for a zoo.
\end{quote}

\subsection*{You'll want to know: \ml{ഉദ്യാനം}}
\textbf{\ml{ഉദ്യാനം}} — \emph{udyānaṁ} — \textquotedblleft{}a public garden; a park\textquotedblright{}.

\begin{quote}
\ml{നഗരത്തിലെ} \ml{ഉദ്യാനത്തിൽ} \ml{വൈകുന്നേരം} \ml{ധാരാളം} \ml{ആളുകൾ} \ml{ഉണ്ടാകും}.
\end{quote}

\begin{quote}
\emph{nagarattile udyānattil vaikunnēraṁ dhārāḷaṁ āḷukaḷ uṇṭākuṁ.}
\end{quote}

\begin{quote}
In the evening there are many people in the city park.
\end{quote}

A national park is a \ml{ദേശീയോദ്യാനം}. A home garden is a \ml{പൂന്തോട്ടം}.

\begin{cousinweb}
\textbf{\ml{ഉദ്യാനം}} (\emph{udyānaṁ}) is borrowed from Sanskrit \emph{udyāna}, \textquotedblleft{}a going out\textquotedblright{} and also \textquotedblleft{}a park, a garden\textquotedblright{} (Monier-Williams).
\end{cousinweb}

\begin{grammarlens}[title={What you've built}]
One more word for getting around.
\end{grammarlens}

\subsection*{Your turn}
\begin{itemize}
\raggedright
  \item \emph{Say it:} \emph{udyānaṁ}
  \item \emph{Say it:} \emph{udyānaṁ}, once more
  \item \emph{Say it:} \emph{nagarattile udyānattil vaikunnēraṁ dhārāḷaṁ āḷukaḷ uṇṭākuṁ.} — In the evening there are many people in the city park.
  \item \emph{Recall:} say the Malayalam for to build a nest, then the Malayalam for a zoo, then say \emph{udyānaṁ} again
\end{itemize}

\begin{tcolorbox}[breakable,colback=teal!4,colframe=teal!35!black,title={Before you move on}]
What does \ml{ഉദ്യാനം} mean? (\textquotedblleft{}A public garden; a park\textquotedblright{}.) Say it once more.
\end{tcolorbox}
\section[\texorpdfstring{\ml{തിമിംഗലം} (timiṁgalaṁ)}{തിമിംഗലം (timiṁgalaṁ)}]{\ml{തിമിംഗലം} (timiṁgalaṁ) — a whale}

\label{lesson:ML-C501-timimgalam}



\begin{quote}
Before the new one: say the Malayalam for a zoo, then the Malayalam for a public garden.
\end{quote}

\subsection*{You'll want to know: \ml{തിമിംഗലം}}
\textbf{\ml{തിമിംഗലം}} — \emph{timiṁgalaṁ} — \textquotedblleft{}a whale\textquotedblright{}.

\begin{quote}
\ml{തിമിംഗലം} \ml{ഒരു} \ml{മീനല്ല}.
\end{quote}

\begin{quote}
\emph{timiṁgalaṁ oru mīnalla.}
\end{quote}

\begin{quote}
A whale is not a fish.
\end{quote}

Plural: \ml{തിമിംഗലങ്ങൾ}, \emph{timiṁgalaṅṅaḷ}.

The biggest animal of the sea; a stranded one is news on the coast.

\begin{cousinweb}
\textbf{\ml{തിമിംഗലം}} (\emph{timiṁgalaṁ}) is borrowed from Sanskrit \emph{timiṅgila}, \textquotedblleft{}swallower of the \emph{timi}\textquotedblright{}, a huge fish of the sea (Monier-Williams).
\end{cousinweb}

\begin{grammarlens}[title={What you've built}]
One more word for getting around.
\end{grammarlens}

\subsection*{Your turn}
\begin{itemize}
\raggedright
  \item \emph{Say it:} \emph{timiṁgalaṁ}
  \item \emph{Say it:} \emph{timiṁgalaṁ}, once more
  \item \emph{Say it:} \emph{timiṁgalaṁ oru mīnalla.} — A whale is not a fish.
  \item \emph{Recall:} say the Malayalam for a zoo, then the Malayalam for a public garden, then say \emph{timiṁgalaṁ} again
\end{itemize}

\begin{tcolorbox}[breakable,colback=teal!4,colframe=teal!35!black,title={Before you move on}]
What does \ml{തിമിംഗലം} mean? (\textquotedblleft{}A whale\textquotedblright{}.) Say it once more.
\end{tcolorbox}
\section[\texorpdfstring{\ml{പ്രവേശനഫീസ്} (pravēśanaphīsŭ)}{പ്രവേശനഫീസ് (pravēśanaphīsŭ)}]{\ml{പ്രവേശനഫീസ്} (pravēśanaphīsŭ) — an entry fee}

\label{lesson:ML-C501-pravesanaphisu}



\begin{quote}
Before the new one: say the Malayalam for a public garden, then the Malayalam for a whale.
\end{quote}

\subsection*{You'll want to know: \ml{പ്രവേശനഫീസ്}}
\textbf{\ml{പ്രവേശനഫീസ്}} — \emph{pravēśanaphīsŭ} — \textquotedblleft{}an entry fee\textquotedblright{}.

\begin{quote}
\ml{മൃഗശാലയിൽ} \ml{കുട്ടികൾക്ക്} \ml{പ്രവേശനഫീസ്} \ml{ഇല്ല}.
\end{quote}

\begin{quote}
\emph{mṛgaśālayil kuṭṭikaḷkkŭ pravēśanaphīsŭ illa.}
\end{quote}

\begin{quote}
There is no entry fee for children at the zoo.
\end{quote}

\ml{പ്രവേശനം}, entry, plus \ml{ഫീസ്}, a fee.

\begin{cousinweb}
\textbf{\ml{പ്രവേശനഫീസ്}} (\emph{pravēśanaphīsŭ}) is a Malayalam compound built from \ml{പ്രവേശനം}, entry, from Sanskrit \emph{praveśana}, \textquotedblleft{}entering\textquotedblright{} (Monier-Williams), and \ml{ഫീസ്}, which was borrowed from English \emph{fees}.
\end{cousinweb}

\begin{grammarlens}[title={What you've built}]
One more word for getting around.
\end{grammarlens}

\subsection*{Your turn}
\begin{itemize}
\raggedright
  \item \emph{Say it:} \emph{pravēśanaphīsŭ}
  \item \emph{Say it:} \emph{pravēśanaphīsŭ}, once more
  \item \emph{Say it:} \emph{mṛgaśālayil kuṭṭikaḷkkŭ pravēśanaphīsŭ illa.} — There is no entry fee for children at the zoo.
  \item \emph{Recall:} say the Malayalam for a public garden, then the Malayalam for a whale, then say \emph{pravēśanaphīsŭ} again
\end{itemize}

\begin{tcolorbox}[breakable,colback=teal!4,colframe=teal!35!black,title={Before you move on}]
What does \ml{പ്രവേശനഫീസ്} mean? (\textquotedblleft{}An entry fee\textquotedblright{}.) Say it once more.
\end{tcolorbox}
\section[\texorpdfstring{\ml{ഒഴുകിയെത്തുക} (oḻukiyettuka)}{ഒഴുകിയെത്തുക (oḻukiyettuka)}]{\ml{ഒഴുകിയെത്തുക} (oḻukiyettuka) — to drift in; to wash in}

\label{lesson:ML-C501-olukiyettuka}



\begin{quote}
Before the new one: say the Malayalam for a whale, then the Malayalam for an entry fee.
\end{quote}

\subsection*{You'll want to know: \ml{ഒഴുകിയെത്തുക}}
\textbf{\ml{ഒഴുകിയെത്തുക}} — \emph{oḻukiyettuka} — \textquotedblleft{}to drift in; to wash in\textquotedblright{}.

\begin{quote}
\ml{കടലിൽ} \ml{നിന്ന്} \ml{പ്ലാസ്റ്റിക്} \ml{കുപ്പികൾ} \ml{തീരത്ത്} \ml{ഒഴുകിയെത്തി}.
\end{quote}

\begin{quote}
\emph{kaṭalil ninnŭ plāsṟṟikŭ kuppikaḷ tīrattŭ oḻukiyetti.}
\end{quote}

\begin{quote}
Plastic bottles drifted in from the sea onto the shore.
\end{quote}

Past: \ml{ഒഴുകിയെത്തി}, \emph{oḻukiyetti}.

\ml{ഒഴുകി}, floating, plus \ml{എത്തുക}, to arrive.

\begin{cousinweb}
\textbf{\ml{ഒഴുകിയെത്തുക}} (\emph{oḻukiyettuka}) is a native Dravidian compound of \ml{ഒഴുകുക}, to flow or float (DEDR 1010), and \ml{എത്തുക}, to reach (DEDR 809).
\end{cousinweb}

\begin{grammarlens}[title={What you've built}]
One more word for getting around.
\end{grammarlens}

\subsection*{Your turn}
\begin{itemize}
\raggedright
  \item \emph{Say it:} \emph{oḻukiyettuka}
  \item \emph{Say it:} \emph{oḻukiyettuka}, once more
  \item \emph{Say it:} \emph{kaṭalil ninnŭ plāsṟṟikŭ kuppikaḷ tīrattŭ oḻukiyetti.} — Plastic bottles drifted in from the sea onto the shore.
  \item \emph{Recall:} say the Malayalam for a whale, then the Malayalam for an entry fee, then say \emph{oḻukiyettuka} again
\end{itemize}

\begin{tcolorbox}[breakable,colback=teal!4,colframe=teal!35!black,title={Before you move on}]
What does \ml{ഒഴുകിയെത്തുക} mean? (\textquotedblleft{}To drift in; to wash in\textquotedblright{}.) Say it once more.
\end{tcolorbox}
